\documentclass[11pt]{article}

% ============================================================
% Packages
% ============================================================
\usepackage[margin=1in]{geometry}
\usepackage{amsmath,amssymb}
\usepackage{array}
\usepackage{booktabs}
\usepackage{enumitem}
\usepackage[T1]{fontenc}
\usepackage{lmodern}
\usepackage{microtype}
\usepackage{graphicx}

% ============================================================
% General Formatting
% ============================================================
\setlength{\parindent}{0pt}
\setlength{\parskip}{0.65em}

\setlist[itemize]{leftmargin=1.5em}
\setlist[enumerate]{leftmargin=1.75em}

\renewcommand{\arraystretch}{1.15}

\graphicspath{{./}}

% ============================================================
% WHITEBOARD Environment
% Everything inside is intended to be written, drawn, or
% displayed as the live whiteboard packet for that topic.
% ============================================================
\newenvironment{whiteboard}
{
    \par\medskip
    \hrule
    \vspace{0.35em}
    \noindent\textbf{\large WHITEBOARD}
    \par\vspace{0.3em}
    \hrule
    \vspace{0.65em}
}
{
    \vspace{0.45em}
    \hrule
    \par\medskip
}

% ============================================================
% Textbook Reference Helpers
% ============================================================
\newcommand{\corresponds}[1]{%
    \textit{Corresponds to Mano, Fig.~#1.}\par\smallskip
}

\newcommand{\seeMano}[1]{%
    \textit{See Mano, Fig.~#1.}\par\smallskip
}

% ============================================================
% Document
% ============================================================
\begin{document}

\begin{center}
    {\LARGE\textbf{CS2614 --- Computer System Architecture}}\\[0.45em]
    {\Large\textbf{Chapter 1: Digital Logic Circuits}}\\[0.45em]
    \textit{Based on M. Morris Mano,
    Computer System Architecture, 3rd ed.}
\end{center}

% ============================================================
\section{Digital Computers and Binary Information}
% ============================================================

A digital computer represents information using variables that take a limited
number of discrete values. In practice, digital computers use two states,
represented by the binary digits 0 and 1. A binary digit is called a
\emph{bit}.

Groups of bits may represent numbers, characters, instructions, or control
information.

\begin{whiteboard}

\[
\boxed{\text{digital}\Rightarrow\text{discrete values}}
\qquad
\boxed{\text{binary}\Rightarrow\{0,1\}}
\]

\[
\boxed{\text{binary digit}=\text{bit}}
\]

\[
\boxed{\text{Information is represented by groups of bits.}}
\]

\end{whiteboard}

For example, when interpreted as a binary number,

\[
1001011_2
=
1(2^6)+1(2^3)+1(2^1)+1(2^0)
=
75_{10}.
\]

The same physical bit pattern may have a different meaning under a different
code.

\begin{whiteboard}

\[
1001011_2=75_{10}
\]

\[
\boxed{
\text{same bits}
\rightarrow
\text{number, character, or control information}
}
\]

\end{whiteboard}

% ============================================================
\section{Hardware, Software, and the Computer System}
% ============================================================

A computer system consists of hardware and software.

Hardware is the physical equipment. Software consists of programs and data.
A sequence of instructions for the computer is called a \emph{program}.

The major hardware portions introduced by Mano are the CPU, memory, and
input-output facilities.

\begin{whiteboard}

\seeMano{1-1}

\[
\boxed{\text{Computer System}=\text{Hardware}+\text{Software}}
\]

\[
\boxed{\text{CPU: manipulate data + fetch/execute instructions}}
\]

\[
\boxed{\text{Memory: store instructions + data}}
\]

\[
\boxed{\text{I/O: communicate with the outside world}}
\]

\end{whiteboard}

Mano distinguishes computer organization, computer design, and computer
architecture.

\begin{whiteboard}

\[
\begin{array}{c|l}
\text{Organization}
&
\text{How hardware operates and connects}
\\
\text{Design}
&
\text{What hardware is used and how it is built}
\\
\text{Architecture}
&
\text{Structure and behavior seen by the user}
\end{array}
\]

\end{whiteboard}

% ============================================================
\section{Logic Gates}
% ============================================================

Binary information is represented physically by two recognizable signal
states.

Logic gates are hardware blocks that perform logical operations on those
binary signals.

A gate can be described by its graphic symbol, Boolean expression, and truth
table.

\begin{whiteboard}

\corresponds{1-2}

\begin{center}
\includegraphics[
    height=0.90\textheight,
    keepaspectratio
]{figure_1_2_DigitalLogicGates.png}
\end{center}

\[
\boxed{
\text{graphic symbol}
\longleftrightarrow
\text{Boolean expression}
\longleftrightarrow
\text{truth table}
}
\]

\end{whiteboard}

The three fundamental operations are AND, OR, and complement.

NAND and NOR are complemented forms of AND and OR. XOR is 1 when its two
inputs differ; XNOR is its complement.

For more than two inputs, Mano describes XOR as an \emph{odd function}: the
output is 1 when an odd number of inputs are 1.

\begin{whiteboard}

\[
\boxed{\text{AND: }x=AB}
\qquad
\boxed{\text{OR: }x=A+B}
\qquad
\boxed{\text{NOT: }x=A'}
\]

\[
\boxed{\text{NAND: }x=(AB)'}
\qquad
\boxed{\text{NOR: }x=(A+B)'}
\]

\[
\boxed{\text{XOR: }x=A\oplus B}
\qquad
\boxed{\text{XNOR: }x=(A\oplus B)'}
\]

\[
A\oplus B=1
\iff
A\neq B
\]

\[
\boxed{\text{XOR is an odd function.}}
\]

\end{whiteboard}

% ============================================================
\section{Boolean Algebra}
% ============================================================

Boolean algebra deals with binary variables and the operations AND, OR, and
complement.

A Boolean function may be represented algebraically, by a truth table, or by
a logic diagram.

For example,

\[
F=x+y'z.
\]

\begin{whiteboard}

\seeMano{1-3}

\[
\boxed{\text{Boolean variables}\in\{0,1\}}
\]

\[
AB\;(\text{AND}),
\qquad
A+B\;(\text{OR}),
\qquad
A'\;(\text{complement})
\]

\[
F=x+y'z
\]

\[
\boxed{
\text{expression}
\longleftrightarrow
\text{truth table}
\longleftrightarrow
\text{logic diagram}
}
\]

\end{whiteboard}

The following identities are especially useful when simplifying Boolean
expressions.

\begin{whiteboard}

\[
\boxed{\text{Useful Boolean Identities}}
\]

\[
x+0=x,
\qquad
x\cdot1=x,
\qquad
x+1=1,
\qquad
x\cdot0=0
\]

\[
x+x=x,
\qquad
xx=x,
\qquad
x+x'=1,
\qquad
xx'=0
\]

\[
x(y+z)=xy+xz
\qquad
(x')'=x
\]

\end{whiteboard}

Simplification matters because a simpler Boolean expression can require fewer
logic gates.

Mano's example is

\[
F=ABC+ABC'+A'C.
\]

Factoring the first two terms gives

\[
F
=
AB(C+C')+A'C
=
AB+A'C.
\]

\begin{whiteboard}

\seeMano{1-6}

\[
F=ABC+ABC'+A'C
\]

\[
=AB(C+C')+A'C
\]

\[
C+C'=1
\]

\[
\boxed{F=AB+A'C}
\]

\[
\boxed{
\text{simpler expression}
\Rightarrow
\text{simpler hardware}
}
\]

\end{whiteboard}

% ============================================================
\section{DeMorgan's Theorems}
% ============================================================

DeMorgan's theorems are

\[
(x+y)'=x'y'
\]

and

\[
(xy)'=x'+y'.
\]

They are important for complementing Boolean expressions and for
understanding the alternate forms of NAND and NOR gates.

\begin{whiteboard}

\[
\boxed{(x+y)'=x'y'}
\qquad
\boxed{(xy)'=x'+y'}
\]

\[
\boxed{
\text{Complement each variable and interchange AND}
\leftrightarrow
\text{OR}
}
\]

\end{whiteboard}

Mano's alternate NOR and NAND symbols are direct applications of
DeMorgan's theorem.

\begin{whiteboard}

\corresponds{1-4 and 1-5}

\begin{center}
\includegraphics[
    width=0.78\textwidth
]{figure_1_4_and_1_5.png}
\end{center}

\[
\boxed{(x+y)'=x'y'}
\qquad
\boxed{(xy)'=x'+y'}
\]

\end{whiteboard}

% ============================================================
\section{Karnaugh Maps}
% ============================================================

Boolean algebra can simplify expressions, but the next algebraic step is not
always obvious.

The Karnaugh map, or K-map, provides a visual procedure for simplification.

Each square represents one minterm. For a function of \(n\) binary variables,
there are \(2^n\) minterms.

\begin{whiteboard}

\[
\boxed{\text{minterm}=\text{one input combination}}
\]

\[
n\text{ variables}
\Rightarrow
2^n\text{ minterms}
\]

\[
ABC=101
\Rightarrow
AB'C
\Rightarrow
m_5
\]

\[
\boxed{
F(A,B,C)=\Sigma(3,4,6,7)
}
\]

\end{whiteboard}

The cells are arranged so that adjacent squares differ in only one variable.

Opposite edges are also adjacent. Groups contain a power-of-two number of
cells.

\begin{whiteboard}

\corresponds{1-7}

\begin{center}
\includegraphics[
    width=0.68\textwidth
]{figure_1_7.png}
\end{center}

\[
\boxed{
\text{Adjacent squares differ in one variable}
}
\]

\[
\boxed{\text{Opposite edges are adjacent}}
\qquad
\boxed{\text{Group sizes: }1,2,4,8,\ldots}
\]

\end{whiteboard}

% ------------------------------------------------------------
\subsection{A Basic Three-Variable Example}
% ------------------------------------------------------------

For Mano's example

\[
F(A,B,C)=\Sigma(3,4,6,7),
\]

the groups produce the terms \(BC\) and \(AC'\).

Therefore,

\[
F=BC+AC'.
\]

\begin{whiteboard}

\corresponds{1-8}

\begin{center}
\includegraphics[
    width=0.50\textwidth
]{figure_1_8.png}
\end{center}

\[
3,7
\Rightarrow
BC
\]

\[
4,6
\Rightarrow
AC'
\]

\[
\boxed{
F=BC+AC'
}
\]

\end{whiteboard}

% ------------------------------------------------------------
\subsection{Standard Forms and Don't-Care Conditions}
% ------------------------------------------------------------

Grouping 1s gives a sum-of-products form.

Grouping 0s can be used to obtain a product-of-sums form.

A don't-care condition may be taken as either 0 or 1 if that choice simplifies
the function.

\begin{whiteboard}

\seeMano{1-11 through 1-13}

\[
\boxed{
\text{Group 1s}
\Rightarrow
\text{sum of products}
}
\]

\[
\boxed{
\text{Group 0s}
\Rightarrow
\text{product of sums}
}
\]

\[
\boxed{
x=
\text{don't care: use as 0 or 1 only if helpful}
}
\]

\end{whiteboard}

For Mano's don't-care example,

\[
F(A,B,C)=\Sigma(0,2,6)
\]

and

\[
d(A,B,C)=\Sigma(1,3,5).
\]

Using the useful don't-care entries gives

\[
F=A'+BC'.
\]

\begin{whiteboard}

\corresponds{1-14}

\begin{center}
\includegraphics[
    width=0.52\textwidth
]{figure_1_14.png}
\end{center}

\[
\boxed{
F=A'+BC'
}
\]

\end{whiteboard}

% ============================================================
\section{Combinational Circuits}
% ============================================================

A combinational circuit is an interconnection of logic gates whose outputs at
any time are functions only of the present inputs.

It contains no stored state.

\begin{whiteboard}

\seeMano{1-15}

\[
\boxed{
\text{Combinational circuit: Output}
=
f(\text{present inputs})
}
\]

\[
\boxed{\text{No stored state}}
\]

\end{whiteboard}

Mano distinguishes analysis from design.

Analysis begins with a circuit and finds its Boolean functions or truth table.

Design begins with a problem statement and ends with a logic diagram.

\begin{whiteboard}

\[
\boxed{
\text{Analysis: circuit}
\rightarrow
\text{Boolean functions}
\rightarrow
\text{truth table}
}
\]

\[
\boxed{
\text{Design: problem}
\rightarrow
\text{variables}
\rightarrow
\text{truth table}
\rightarrow
\text{simplify}
\rightarrow
\text{logic diagram}
}
\]

\end{whiteboard}

% ============================================================
\section{Half-Adder and Full-Adder}
% ============================================================

A half-adder adds two binary digits \(x\) and \(y\).

It produces a sum \(S\) and carry \(C\).

\[
S=x\oplus y,
\qquad
C=xy.
\]

\begin{whiteboard}

\seeMano{1-16}

\[
\boxed{\text{Half-Adder}}
\]

\[
\begin{array}{cc|cc}
x&y&C&S\\
\hline
0&0&0&0\\
0&1&0&1\\
1&0&0&1\\
1&1&1&0
\end{array}
\]

\[
\boxed{S=x\oplus y}
\qquad
\boxed{C=xy}
\]

\end{whiteboard}

A full-adder adds two significant bits and a carry from the previous lower
position.

Mano uses inputs \(x\), \(y\), and \(z\), with outputs \(S\) and \(C\).

\[
S=x\oplus y\oplus z
\]

and

\[
C=xy+(x\oplus y)z.
\]

The circuit can be constructed from two half-adders and an OR gate.

\begin{whiteboard}

\corresponds{1-18}

\begin{center}
\includegraphics[
    width=0.78\textwidth
]{figure_1_18.png}
\end{center}

\[
\boxed{
S=x\oplus y\oplus z
}
\]

\[
\boxed{
C=xy+(x\oplus y)z
}
\]

\[
\boxed{
\text{Full-adder}
=
\text{two half-adders}
+
\text{OR}
}
\]

\end{whiteboard}

% ============================================================
\section{Sequential Circuits and Flip-Flops}
% ============================================================

A sequential circuit contains storage.

Its behavior therefore depends on both the present input and the present state
of the storage elements.

\begin{whiteboard}

\[
\boxed{
\text{Combinational: }Y=f(X)
}
\]

\[
\boxed{
\text{Sequential: }Y=f(X,\text{present state})
}
\]

\[
\boxed{
\text{next state}
=
g(X,\text{present state})
}
\]

\end{whiteboard}

Clocked synchronous sequential circuits use a clock signal to coordinate state
changes.

A flip-flop is a binary storage element capable of storing one bit.

\begin{whiteboard}

\[
\boxed{
\text{Clock}
\Rightarrow
\text{synchronizes state changes}
}
\]

\[
\boxed{
\text{Flip-flop}
=
\text{1-bit storage element}
}
\]

\[
Q(t)=\text{present state}
\qquad
Q(t+1)=\text{next state}
\]

\end{whiteboard}

% ------------------------------------------------------------
\subsection{SR Flip-Flop}
% ------------------------------------------------------------

The SR flip-flop provides set and reset inputs.

The input combination

\[
S=R=1
\]

produces an indeterminate next state.

\begin{whiteboard}

\corresponds{1-19}

\begin{center}
\includegraphics[
    width=0.66\textwidth
]{figure_1_19.png}
\end{center}

\[
\boxed{
S=R=1
\Rightarrow
\text{indeterminate}
}
\]

\end{whiteboard}

% ------------------------------------------------------------
\subsection{D Flip-Flop}
% ------------------------------------------------------------

For a D flip-flop, the next state is simply the D input sampled at the active
clock transition:

\[
Q(t+1)=D.
\]

\begin{whiteboard}

\corresponds{1-20}

\begin{center}
\includegraphics[
    width=0.62\textwidth
]{figure_1_20.png}
\end{center}

\[
\boxed{
Q(t+1)=D
}
\]

\[
\boxed{
\text{active clock transition: }Q\leftarrow D
}
\]

\end{whiteboard}

% ------------------------------------------------------------
\subsection{JK and T Flip-Flops}
% ------------------------------------------------------------

The JK flip-flop replaces the indeterminate SR condition with a toggle.

When

\[
J=K=1,
\]

the next state is the complement of the present state.

The T flip-flop has a single toggle input.

When \(T=0\), the state does not change. When \(T=1\), the state complements.

\begin{whiteboard}

\corresponds{1-21 and 1-22}

\begin{center}
\includegraphics[
    width=0.67\textwidth
]{figure_1_21_and_1_22.png}
\end{center}

\[
\boxed{
J=K=1
\Rightarrow
\text{toggle}
}
\]

\[
\boxed{
Q(t+1)=Q(t)\oplus T
}
\]

\end{whiteboard}

% ------------------------------------------------------------
\subsection{Edge Triggering}
% ------------------------------------------------------------

An edge-triggered flip-flop changes state only at a designated transition of
the clock.

A positive-edge device responds to

\[
0\rightarrow1,
\]

while a negative-edge device responds to

\[
1\rightarrow0.
\]

The data input must also satisfy setup-time and hold-time requirements around
the active transition.

\begin{whiteboard}

\corresponds{1-23}

\begin{center}
\includegraphics[
    width=0.78\textwidth
]{figure_1_23.png}
\end{center}

\[
\boxed{
\text{positive edge: }0\rightarrow1
}
\qquad
\boxed{
\text{negative edge: }1\rightarrow0
}
\]

\[
\boxed{
\text{setup time: input stable before the active edge}
}
\]

\[
\boxed{
\text{hold time: input stable after the active edge}
}
\]

\end{whiteboard}

% ============================================================
\section{Characteristic and Excitation Tables}
% ============================================================

A characteristic table specifies the next state when the flip-flop inputs and
present state are known.

An excitation table works in the reverse direction: it specifies which inputs
are needed to obtain a desired state transition.

\begin{whiteboard}

\[
\boxed{
\text{Characteristic: inputs}+Q(t)
\rightarrow
Q(t+1)
}
\]

\[
\boxed{
\text{Excitation: }Q(t)+Q(t+1)
\rightarrow
\text{required FF inputs}
}
\]

\end{whiteboard}

The JK excitation table is especially useful in Mano's counter design.

\begin{whiteboard}

\[
\boxed{\text{JK Excitation Table}}
\]

\[
\begin{array}{cc|cc}
Q(t)&Q(t+1)&J&K\\
\hline
0&0&0&x\\
0&1&1&x\\
1&0&x&1\\
1&1&x&0
\end{array}
\]

\[
\boxed{x=\text{don't care}}
\]

\end{whiteboard}

% ============================================================
\section{Sequential-Circuit Description and Design}
% ============================================================

A clocked sequential circuit combines combinational logic and flip-flops.

The combinational part receives external inputs and present-state values. It
produces external outputs and the signals that determine the next flip-flop
states.

Mano's example with two D flip-flops uses the input equations

\[
D_A=Ax+Bx,
\qquad
D_B=A'x,
\]

and the external output

\[
y=Ax'+Bx'.
\]

\begin{whiteboard}

\seeMano{1-24 and 1-25}

\[
\boxed{
\text{next state}
=
f(\text{present state, inputs})
}
\]

\[
\boxed{
\text{outputs}
=
g(\text{present state, inputs})
}
\]

\[
\boxed{D_A=Ax+Bx}
\qquad
\boxed{D_B=A'x}
\]

\[
\boxed{y=Ax'+Bx'}
\]

\end{whiteboard}

A state table relates present state and input to next state and output.

A state diagram represents the same information graphically.

\begin{whiteboard}

\seeMano{1-26}

\[
\boxed{
\text{present state + input}
\rightarrow
\text{next state + output}
}
\]

\[
\boxed{
\text{state diagram: circle}
=
\text{state}
}
\]

\[
\boxed{
\text{arrow}
=
\text{transition}
}
\qquad
\boxed{
x/y
=
\text{input/output}
}
\]

\end{whiteboard}

The general design path is:

\begin{whiteboard}

\[
\boxed{
\text{specification}
\rightarrow
\text{state diagram}
\rightarrow
\text{state table}
}
\]

\[
\boxed{
\rightarrow
\text{choose FF type}
\rightarrow
\text{excitation table}
\rightarrow
\text{K-maps}
}
\]

\[
\boxed{
\rightarrow
\text{FF input equations}
\rightarrow
\text{logic circuit}
}
\]

\end{whiteboard}

% ============================================================
\section{Design Example: 2-Bit Binary Counter}
% ============================================================

Mano's final example is a 2-bit binary counter controlled by an input \(x\).

When \(x=1\), the circuit advances through the binary count sequence.

When \(x=0\), the present state is held.

\begin{whiteboard}

\corresponds{1-27}

\begin{center}
\includegraphics[
    width=0.48\textwidth
]{figure_1_27.png}
\end{center}

\[
\boxed{
x=1:
\quad
00
\rightarrow
01
\rightarrow
10
\rightarrow
11
\rightarrow
00
}
\]

\[
\boxed{
x=0:
\quad
\text{hold present state}
}
\]

\end{whiteboard}

Using the JK excitation table and K-map simplification, Mano obtains the four
flip-flop input equations

\[
J_A=Bx,
\qquad
K_A=Bx,
\qquad
J_B=x,
\qquad
K_B=x.
\]

\begin{whiteboard}

\corresponds{1-28}

\begin{center}
\includegraphics[
    width=0.70\textwidth
]{figure_1_28.png}
\end{center}

\[
\boxed{
J_A=K_A=Bx
}
\]

\[
\boxed{
J_B=K_B=x
}
\]

\end{whiteboard}

The resulting hardware contains two JK flip-flops and one AND gate.

\begin{whiteboard}

\corresponds{1-29}

\begin{center}
\includegraphics[
    width=0.54\textwidth
]{figure_1_29.png}
\end{center}

\[
\boxed{
\text{state diagram}
\rightarrow
\text{state table}
\rightarrow
\text{excitation table}
}
\]

\[
\boxed{
\rightarrow
\text{K-maps}
\rightarrow
\text{FF equations}
\rightarrow
\text{hardware}
}
\]

\end{whiteboard}

% ============================================================
\section{Chapter Summary}
% ============================================================

Chapter 1 begins with binary information and logic gates, develops Boolean
algebra and K-map simplification, uses them to design combinational circuits,
and then introduces flip-flops and stored state to build sequential circuits.

\begin{whiteboard}

\[
\boxed{
\text{bits}
\rightarrow
\text{logic gates}
\rightarrow
\text{Boolean algebra}
\rightarrow
\text{K-maps}
}
\]

\[
\boxed{
\rightarrow
\text{combinational circuits}
\rightarrow
\text{flip-flops}
\rightarrow
\text{sequential circuits}
}
\]

\[
\boxed{
\text{Combinational: present input only}
}
\]

\[
\boxed{
\text{Sequential: present input + present state}
}
\]

\end{whiteboard}

\end{document}